\chapter{State of the Art}

Existing approaches for modelling and predicting driver behaviour can be broadly grouped into rule-based models, probabilistic models, and learning-based data-driven models. 
In addition, hybrid methods combine elements from more than one of these families in order to exploit their complementary strengths. 


\section{Rule-Based Models}

Rule-based models characterize driving behaviour through explicit mathematical or logical rules. 
These models relate observable variables such as spacing, relative velocity, and desired speed to longitudinal acceleration or lane-change decisions. 
Their main advantages are interpretability and low computational cost \cite{zhou2025}. For this reason, they are widely used in traffic simulation, 
\ac{ADAS}, and baseline behaviour modelling.  
However, as they rely on hand-crafted assumptions about driver decision-making, they often represent behaviour in a deterministic or simplified manner. 


\subsection{Car-Following Models}

\subsubsection{Intelligent Driver Model}
The \ac{IDM} is one of the most widely used car-following models \cite{zhou2025}. 
It expresses a vehicle's acceleration as a continuous function of its current velocity, 
desired velocity, spacing to the leading vehicle, and relative velocity \cite{PhysRevE.62.1805}. The formulation also includes parameters related to desired time headway and comfortable 
braking, which make the model suitable for safety-oriented traffic simulation. 
Because of its simple structure and physically meaningful parameters, the \ac{IDM} remains a common reference model in intelligent driving research.

\subsubsection{Gipps' Model}
Gipps' model is one of the earliest rule-based car-following models. 
It is based on the idea that a driver chooses a speed that is both desirable and safe with respect to the 
preceding vehicle \cite{GIPPS1981105}. This safety-oriented formulation
allows the model to capture conservative driving behaviour and ensures collision-free motion under typical
traffic conditions.

\subsubsection{Optimal Velocity Model and Full Velocity Difference Model}
Another classical approach is the \ac{OVM}, in which the driver's acceleration depends on the difference 
between the current velocity and an optimal velocity determined by the spacing to the leading vehicle \cite{PhysRevE.51.1035}. This model captures
the driver's intuition to accelerate when the gap ahead increases and decelerate when it decreases. An extension of the \ac{OVM} is the \ac{FVDM}, 
where, instead of only considering the gap between the vehicles, the model also captures
the velocity difference between the ego and preceding vehicles when calculating the acceleration \cite{PhysRevE.64.017101}. This
helps reduce the oscillations observed in the original \ac{OVM}.


\subsection{Lane-Change Models}

\subsubsection{MOBIL Model}
The \ac{MOBIL} model is a lane-changing model built upon the \ac{IDM} to describe lateral driving behaviour \cite{zhou2025} \cite{Kesting2007GeneralLM}. 
It incorporates incentive and safety criteria
to evaluate the feasibility of lane changes. The incentive criterion indicates whether a lane
change could be advantageous in terms of the driver's acceleration, whereas the safety criterion ensures
that the lane change does not cause unnecessary deceleration for surrounding vehicles.


\par Despite their interpretability and efficiency, rule-based models have limited ability to represent uncertainty and inter-driver variability in complex traffic scenarios. 
These limitations motivate the use of probabilistic and learning-based approaches.